% Part: proof-theory
% Chapter: sequent-calculus
% Section: rules-proofs

\documentclass[../../../include/open-logic-section]{subfiles}

\begin{document}

\olfileid{pt}{seq}{rp}

\olsection{规则与\usetoken{P}{proof}}

我们首先给出一些定义并规定一些术语。

例如，考虑 \Log{G3c} 中允许从包含 $!A$ 和~$!B$ 的相继式推出包含 $!A \land !B$ 的相继式的两条规则：
\[
\Axiom$ !A, !B, \Gamma \fCenter \Delta$
\RightLabel{\LeftR{\land}}
\UnaryInf$ !A \land !B, \Gamma \fCenter \Delta$
\DisplayProof
\qquad
\Axiom$\Gamma \fCenter \Delta, !A$
\Axiom$ \Gamma \fCenter \Delta, !B$
\RightLabel{\RightR{\land}}
\BinaryInf$ \Gamma \fCenter \Delta, !A \land !B $
\DisplayProof
\]
横线上方的相继式称为\emph{前提}，横线下方的相继式称为\emph{结论}。每条规则中 $\Gamma$ 和 $\Delta$ 里的公式称为\emph{上下文}或\emph{边公式}。结论中不属于上下文的那一个公式称为\emph{主公式}，前提中不属于上下文的公式称为\emph{活动}（或\emph{辅助}）公式。

量词规则稍复杂一些。例如，\Log{G1c} 中关于 $\lforall$ 的规则是：
\[\Axiom$ !A(t), \Gamma \fCenter \Delta$
\RightLabel{\LeftR{\lforall}}
\UnaryInf$ \lforall[x][!A(x)],\Gamma \fCenter \Delta$
\DisplayProof
\qquad
\Axiom$ \Gamma \fCenter \Delta, !A(a) $
\RightLabel{\RightR{\lforall}}
\UnaryInf$ \Gamma \fCenter \Delta, \lforall[x][!A(x)]$
\DisplayProof
\]
在 \LeftR{\lforall} 规则中，活动公式是~$!A(t)$，其中 $t$~是闭项，即不含变元的项。这种推理是正确的——也就是说，它符合模式规则 \LeftR{\lforall}——如果存在某个公式~$!A$、某个变元~$x$ 和某个闭项~$t$，使得结论中的主公式是 $\lforall[x][!A]$，而前提中的活动公式是~$\Subst{!A}{t}{x}$。\RightR{\lforall} 规则与此类似，区别在于前提中的活动公式不能是任意项~$t$，而必须是~$\Subst{!A}{a}{x}$，其中 $a$ 是一个常项符号。这个常项符号称为该 \RightR{\lforall} 推理的\emph{特征变元}。只有当下方的相继式不包含~$a$ 时——也就是说，$a$ 不在 $\Gamma$、$\Delta$ 或~$!A$ 中出现——该推理才是正确的。这个条件称为\emph{特征变元条件}。

像 \Log{G1c} 和 \Log{G3c} 这样的相继式系统，由一组按模式规定的规则，连同一些（同样按模式规定的）作为公理的相继式构成。这类系统中的推导是从公理出发，使用推理规则归纳生成的相继式的有限树。每个叶节点是一个公理，而所有其他相继式都是通过系统的某条推理规则从其紧上方的相继式推出的。

\begin{defn}\ollabel{defn:proof}
相继式系统中的\emph{推导}是按如下方式归纳定义的相继式的有限树：
\begin{enumerate}
  \item\ollabel{defn:prf:axiom} 每个公理都是一个推导。
  \item\ollabel{defn:prf:one-prem} 如果 $\pi_1$ 是一个以相继式~$S_1$ 结尾的推导，而 $S$ 是由~$S_1$ 通过某条单前提规则~$R$ 推出的相继式，那么
  \[
  \AxiomC{}
  \RightLabel{$\pi_1$}
  \DeduceC{$S_1$}
  \RightLabel{$R$}
  \UnaryInfC{$S$}
  \DisplayProof
  \]
  是一个推导。
  \item\ollabel{defn:prf:two-prem} 如果 $\pi_1$ 和 $\pi_2$ 分别是以相继式~$S_1$ 和~$S_2$ 结尾的推导，并且 $S$ 是由~$S_1$ 和$S_2$ 通过某条双前提规则~$R$ 推出的相继式，那么
  \[
  \AxiomC{}
  \RightLabel{$\pi_1$}
  \DeduceC{$S_1$}
  \AxiomC{}
  \RightLabel{$\pi_2$}
  \DeduceC{$S_2$}
  \RightLabel{$R$}
  \BinaryInfC{$S$}
  \DisplayProof
  \]
  是一个推导。
\end{enumerate}
\end{defn}

\Olref{defn:proof} 将推导定义为相继式的树。我们设想这些树是“向上”生长的。这棵树最顶端的（叶）节点是公理，称为\emph{初始相继式}。最底部的相继式称为\emph{末端相继式}。如果 $\pi$ 是一个以相继式~$S$ 为末端相继式的推导，我们也写作 $\pi \Proves S$。如果存在相继式~$S$ 的一个推导，我们写作 $\Proves S$，也可以在 $\Proves$~符号左边标明证明所在的系统，例如 $\Log{G1c} \Proves !A, !B \Sequent !A \land
!B$。推导中除初始相继式外的每一个相继式，都是通过某条规则~$R$ 所做推理的\emph{结论}。在推导中紧挨着一个相继式上方的一个或两个相继式（即该节点的父节点）称为该推理的\emph{前提}。

在归纳构造一个推导时所用到的推导称为该推导的\emph{子推导}。那些以某推理的前提为末端相继式的子推导，称为该推理的\emph{直接子推导}。

常常需要用一个自然数来度量推导的复杂度。我们可以简单地计算推导中相继式的个数，但通常更有用的度量是推导的\emph{高度}：它是末端相继式与某个初始相继式之间的相继式的最大数目。我们也可以给出一个归纳定义。

\begin{defn}
推导~$\pi$ 的\emph{高度}~$\pheight{\pi}$ 按如下方式归纳定义。
\begin{enumerate}
  \item 如果 $\pi$ 仅由一个公理组成，则 $\pheight{\pi} = 0$。
  \item 如果 $\pi$ 以一个单前提的推理结束，其直接子推导为~$\pi_1$，则 $\pheight{\pi} = \pheight{\pi_1} + 1$。
  \item 如果 $\pi$ 以一个双前提的推理结束，其直接子推导为~$\pi_1$ 和~$\pi_2$，则 $\pheight{\pi} =
  \max(\pheight{\pi_1}, \pheight{\pi_2}) + 1$。
\end{enumerate}
如果相继式~$S$ 有一个高度 $\le n$ 的推导，我们记作 $\Proves[n] S$。
\end{defn}

\begin{ex}
在 \Log{G3c} 中，任何形如 $!C, \Gamma \Sequent \Delta, !C$ 的相继式都是一条公理，其中 $!C$ 必须是原子的。
假设 $!D$ 和 $!E$ 是原子语句。
那么 $!D, !E \Sequent !D$ 和 $!D, !E \Sequent !E$ 就是公理 $!C, \Gamma \Sequent \Delta, !C$ 的实例。
例如，若我们在 $!C, \Gamma \Sequent \Delta, !C$ 中取 $!C \ident !E$、$\Gamma = \{!D\}$ 且 $\Delta = \emptyset$，得到的正是相继式 $!D, !E \Sequent !E$。（请记住我们处理的是多重集，因此 $!D, !E \Sequent !E$ 和 $!E, !D \Sequent !D$ 是同一个相继式。）

由于 $!D, !E \Sequent !D$ 是 \Log{G3c} 的一条公理的实例，它（本身）就算作 \Log{G3c} 中的一个推导；根据 \olref{defn:proof}\olref{defn:prf:axiom}，$!D, !E \Sequent !E$ 也同样是一个推导。
让我们称它们为 $\pi_1$ 和~$\pi_2$。
根据 \olref{defn:prf:two-prem}，我们可以利用双前提规则~$\RightR{\land}$，从这两个推导生成一个新的推导（$\pi_3$）：
\[
\Axiom$!D, !E \fCenter !D$
\Axiom$!D, !E \fCenter !E$
\RightLabel{\RightR{\land}}
\BinaryInf$!D, !E  \fCenter !D \land !E $
\DisplayProof
\]
接着，我们从 $\pi_3$ 出发，利用单前提规则~$\LeftR{\land}$（使用 \olref{defn:proof}\olref{defn:prf:one-prem}）得到推导~$\pi_4$：
\[
\Axiom$!D, !E \fCenter !D$
\Axiom$!D, !E \fCenter !E$
\RightLabel{\RightR{\land}}
\insertBetweenHyps{\hspace{5ex}}
\BinaryInf$!D, !E  \fCenter !D \land !E $
\LeftSubproofLabel{$\pi_3$}
\RightLabel{\LeftR{\land}}
\UnaryInf$!D \land !E  \fCenter !D \land !E$
\RightSubproofLabel{$\pi_4$}
\DisplayProof
\]
$\pi_4$ 中最后一个推理的直接子推导是~$\pi_3$。
$\RightR{\land}$ 推理的直接子推导是 $\pi_1$ 和~$\pi_2$。
$\pi_4$ 的所有子推导包括 $\pi_1$、$\pi_2$、$\pi_3$ 和 $\pi_4$ 自身。
我们有 $\pheight{\pi_1} =
\pheight{\pi_2} = 0$，$\pheight{\pi_3} = 1$ 且 $\pheight{\pi_4} = 2$。
我们已表明，对于任意的原子语句 $!D$ 和~$!E$，$\Log{G3c} \Proves[4] !D \land !E  \fCenter !D
\land !E$。
\end{ex}

\end{document}